\subsection{Document Layout Analysis}

A page is not a bag of words. The same number may be a table entry, a footnote marker or a page number; its role depends on position and relations to other content. Layout analysis makes those relations available to recognition, retrieval and reasoning. Three outputs should be distinguished: \emph{regions} locate content, \emph{relations} connect it, and \emph{representations} encode both for a downstream model.

\subsubsection{From Regions to Document Structure}

A region detector predicts boxes or masks for text, titles, tables and figures. This provides useful processing units: a recognizer can receive a text crop, a table parser can receive a table region, and a retriever can retain the originating page. Detection is therefore an interface between page pixels and specialist components, rather than a complete account of document meaning.

The main design choice is how much context the detector sees. Local appearance helps distinguish a table from prose, but can be ambiguous when captions, columns and figures are tightly packed. VSR combines visual, semantic and relational features~\cite{zhang2021vsr}; DocLayout-YOLO uses synthetic document data and global-to-local perception~\cite{doclayout_yolo_2024}. These approaches illustrate complementary ways to use context without assuming that all document types share the same layout.

Boxes alone leave important structure unresolved. A two-column page needs reading order; a table needs header--cell links; a figure needs its caption. Recovering these relations can be more consequential for QA than a small improvement in box overlap. For example, reading across columns may create a fluent but nonexistent sentence even when every word is recognized correctly.

Consider a report table whose unit appears above the box and whose qualification appears in a footnote below it. Cropping the cells accurately still removes context needed to interpret their values. A useful structural output therefore includes the table's neighborhood and links to its notes. The region boundary and the boundary of sufficient evidence are not necessarily the same.

\subsubsection{Encoding Text and Spatial Relations}

Layout-aware models pass OCR text and geometry into a learned representation. LayoutLM adds two-dimensional token positions to Transformer inputs~\cite{xu2019layoutlm}. Compared with plain text serialization, the model retains information about which words occur on the same line or near the same field. It still depends on the recognizer: an omitted word is absent from the token input, and an incorrect box changes the spatial evidence.

The representation can separate linguistic and geometric processing. LiLT uses a language-independent layout Transformer that interacts with text features~\cite{wang2022lilt}, making layout transfer a different problem from learning a particular language. This is useful when document templates recur across languages, although script-specific recognition quality remains a separate dependency.

Reading order deserves explicit treatment. XYLayoutLM addresses ordering noise through an augmented XY-cut procedure and positional encoding~\cite{gu2022xylayoutlm}; relation-based work instead represents order as directed links~\cite{zhang2024reading_order_relations}. A sequence is convenient for a language model, while a relation graph can preserve alternatives and hierarchy. The appropriate interface depends on whether the question requires a passage, a table neighborhood or a cross-page link.

\subsubsection{Using Structure During Reasoning}

Structure can condition generation directly. DocLLM uses text and bounding boxes with disentangled attention, without an image encoder~\cite{wang2024docllm}. LayoutLLM introduces layout instruction tuning~\cite{luo2024layoutllm}. Both make geometry available to the answering process; neither should be described as observing all visual content simply because spatial information is present.

Structure can also organize evidence access. LAD-RAG uses layout graphs for cross-page retrieval~\cite{ladrag2025}. Here the benefit is not a new text recognizer but a different retrieval unit: a table, its caption and a linked paragraph can remain connected rather than become unrelated chunks. This is particularly relevant when a question needs both a numerical cell and an explanation elsewhere in the document.

These choices expose different failure points. Coordinate-aware generation can fail after a bad OCR pass; graph retrieval can fail when a required link was never built. The useful question is therefore which relations survive each interface. Real-document analyses show that performance of pretrained text--layout models can degrade outside their familiar distributions~\cite{zhang2024unveiling}; visual plausibility of the detected page is insufficient evidence of robust reasoning.

\subsubsection{What Layout Benchmarks Measure}

Dataset construction shapes the task. PubLayNet obtains large-scale region supervision by aligning scientific PDFs with XML~\cite{zhong2019publaynet}; DocBank provides token-level document annotations~\cite{li2020docbank}; DocLayNet uses human-annotated pages from several document domains~\cite{pfitzmann2022doclaynet}. These resources offer different balances of scale, label granularity and domain diversity.

Region overlap and category accuracy measure detection, while reading-order and relation annotations test structure. Downstream QA adds another requirement: the selected regions must jointly support an answer. Evaluations should connect these levels by inspecting whether an error arose in the box, its semantic role, its links or the answer that used it. A high detection score cannot resolve those questions by itself.
